\chapter{La campagne avec XFOIL}\label{ch:campagne}
\iterbanner{V3}

\section{Protocole}
Exactement celui de la seconde campagne CST du lot L6 (\code{contraintes/experiences\_cst.py}) : filtre géométrique avant simulation, traînée de plus de 200 counts comptée comme un échec, budget de 300 simulations, 8 graines, \emph{mêmes profils initiaux par graine}, ce qui permet une comparaison appariée. Les méthodes de la seconde campagne ne sont pas relancées : leurs exécutions sont relues, et le front de référence est recalculé sur l'union des deux campagnes. Deux méthodes nouvelles, toutes deux avec le modèle en $z$ et le recuit simulé, à 30 qubits :
\begin{itemize}
\item \code{quadz\_lagrange} : $C_{L,\max}$ et les six contraintes de non-croisement en termes lagrangiens, 6 résolutions par QUBO, lectures vérifiées ;
\item \code{quadz\_lagrange\_clf} : la même chose, plus le classifieur de convergence d'XFOIL avec $p_{\min}=0{,}1$.
\end{itemize}
La pénalité déséquilibrée, moins bonne que le lagrangien sur les deux études hors ligne, n'a pas été relancée.

\section{Résultats}
Hypervolume des profils faisables simulés, rapporté à celui du front de référence (11 profils, sur l'union des 20\,602 profils simulés par les deux campagnes), médiane sur 8 graines. « vs actuel » : nombre de graines où la méthode fait mieux que l'approche actuelle à 300 simulations, et $p$ du test de Wilcoxon apparié. À 100 simulations, le classifieur fait mieux sur 5 graines ($p=0{,}84$).
\begin{center}
\renewcommand{\arraystretch}{1.2}
\setlength{\tabcolsep}{4pt}
\small
\begin{tabular}{lcccccc}
\toprule
& à 100 & à 300 & vs actuel (300) & faisables & échecs XFOIL & $C_{L,\max}$ faible \\
\midrule
bits d'écart (actuel) & 0,768 & \textbf{0,882} & --- & 43~\% & 28~\% & 29~\% \\
lagrangien vérifié & 0,768 & 0,839 & 2/8, $p=0{,}38$ & 37~\% & 40~\% & 22~\% \\
lagrangien + classifieur & \textbf{0,815} & 0,840 & 1/8, $p=0{,}15$ & 41~\% & 34~\% & 25~\% \\
\midrule
factorisation (L6) & 0,833 & 0,895 & 6/8, $p=0{,}38$ & 47~\% & 25~\% & 28~\% \\
NSGA-II & 0,718 & 0,813 & 2/8, $p=0{,}04$ & 72~\% & 13~\% & 15~\% \\
hasard & 0,683 & 0,772 & 1/8, $p=0{,}02$ & 76~\% & 6~\% & 17~\% \\
\bottomrule
\end{tabular}
\end{center}

\begin{figure}[htbp]
\centering
% figure générée par tools/make_figures.py
\begin{tikzpicture}
\begin{axis}[width=0.92\linewidth,height=7cm,xmin=0,xmax=300,ymin=0.3,ymax=1,xlabel={simulations},
  ylabel={$HV/HV_{\rm ref}$ (médiane sur 8 graines)},grid=major,grid style={cmuted!20},legend pos=south east,
  legend cell align=left,legend style={font=\footnotesize}]
\addplot[cmuted,thick,dashed,mark size=1.4pt] coordinates {(20,0.61544) (40,0.64622) (60,0.65565) (80,0.66783) (100,0.68305) (120,0.6884) (140,0.72706) (160,0.72737) (180,0.72737) (200,0.72913) (220,0.76737) (240,0.77032) (260,0.77032) (280,0.77082) (300,0.77155)};
\addlegendentry{hasard}
\addplot[cink,thick,dotted,mark size=1.4pt] coordinates {(20,0.61096) (40,0.62004) (60,0.6257) (80,0.65257) (100,0.71827) (120,0.72088) (140,0.74279) (160,0.76578) (180,0.77044) (200,0.77275) (220,0.78817) (240,0.78883) (260,0.789) (280,0.81253) (300,0.81278)};
\addlegendentry{NSGA-II}
\addplot[cblue,very thick,mark=*,mark size=1.4pt] coordinates {(20,0.61204) (40,0.69689) (60,0.72471) (80,0.76812) (100,0.76812) (120,0.76812) (140,0.79083) (160,0.80248) (180,0.82308) (200,0.82499) (220,0.82879) (240,0.84877) (260,0.85732) (280,0.86329) (300,0.88181)};
\addlegendentry{bits d'écart (actuel)}
\addplot[corange,thick,mark=square*,mark size=1.4pt] coordinates {(20,0.63019) (40,0.6879) (60,0.76189) (80,0.8174) (100,0.83258) (120,0.85689) (140,0.86559) (160,0.86807) (180,0.87256) (200,0.87577) (220,0.87915) (240,0.8796) (260,0.88328) (280,0.89162) (300,0.89483)};
\addlegendentry{machine de factorisation (bits d'écart)}
\addplot[cgreen,very thick,mark=triangle*,mark size=1.4pt] coordinates {(20,0.58278) (40,0.62809) (60,0.708) (80,0.75695) (100,0.7682) (120,0.79466) (140,0.79466) (160,0.79466) (180,0.79466) (200,0.79986) (220,0.83212) (240,0.83255) (260,0.83795) (280,0.83895) (300,0.83944)};
\addlegendentry{lagrangien vérifié}
\addplot[cviolet,very thick,mark=diamond*,mark size=1.4pt] coordinates {(20,0.58278) (40,0.69995) (60,0.75154) (80,0.80242) (100,0.81519) (120,0.82286) (140,0.82286) (160,0.82286) (180,0.82286) (200,0.82286) (220,0.82644) (240,0.82822) (260,0.836) (280,0.83936) (300,0.83991)};
\addlegendentry{lagrangien + classifieur}
\end{axis}
\end{tikzpicture}

\caption{Campagne CST : hypervolume relatif en fonction du nombre de simulations (médiane sur 8 graines). Les deux méthodes lagrangiennes (vert, violet) partent comme l'approche actuelle, puis plafonnent.}\label{fig:campagne-hv}
\end{figure}

\paragraph{1. Les études hors ligne ne se confirment pas.} Les deux méthodes lagrangiennes finissent \emph{sous} l'approche actuelle (0,84 contre 0,88), sans que l'écart soit significatif avec 8 graines. Le classifieur aide en début de budget (0,815 à 100 simulations, devant l'approche actuelle sur 5 graines), puis plafonne aussi.

\paragraph{2. Ce qui était prévu se produit.} Les $C_{L,\max}$ trop faibles baissent, de 29~\% à 22~\% : la formulation fait mieux respecter la contrainte du modèle, comme l'annonçait le rejeu (chapitre~\ref{ch:hors-ligne}). Le classifieur fait baisser les échecs d'XFOIL de 40~\% à 34~\% par rapport au lagrangien seul.

\paragraph{3. Ce qui n'était pas prévu : les échecs d'XFOIL montent.} De 28~\% à 40~\% avec le lagrangien seul. La faisabilité totale baisse donc (37~\% contre 43~\%). On a cherché pourquoi dans les profils proposés : la part des coefficients placés à un niveau \emph{extrême} de la grille (0 ou 7).
\begin{center}
\begin{tabular}{lc}
\toprule
& coefficients aux niveaux extrêmes \\
\midrule
hasard & 25~\% \\
machine de factorisation (régularisée) & 30~\% \\
bits d'écart (actuel) & 56~\% \\
lagrangien + classifieur & 62~\% \\
lagrangien vérifié & \textbf{66~\%} \\
\bottomrule
\end{tabular}
\end{center}
Pour les trois méthodes à modèle en $z$, les profils dont au moins 6 coefficients sur 10 sont extrêmes échouent dans XFOIL 40 à 48~\% du temps, contre 16 à 21~\% pour les autres.

\paragraph{4. Le mécanisme.} Le modèle en $z$ est une parabole à 66 coefficients, sans régularisation : son minimum est souvent sur un bord de la grille, là où il extrapole. Avec les bits d'écart, la pénalité occupait 99~\% du spectre ; le recuit, réglé sur l'objectif, l'optimisait mal, et ses lectures restaient plus centrées. \emph{La pénalité régularisait par accident.} Le lagrangien rend tout le spectre à l'objectif : le recuit trouve beaucoup mieux le minimum du modèle, c'est-à-dire ses erreurs d'extrapolation dans les coins, où XFOIL décroche. La machine de factorisation, régularisée, ne va aux extrêmes que 30~\% du temps : c'est elle qu'il faudrait associer au lagrangien.

\paragraph{5. Pourquoi le rejeu ne l'a pas vu.} Il mesurait la faisabilité \emph{prédite} et l'objectif \emph{prédit} du candidat, sur des états produits par l'approche actuelle. Or le défaut du lagrangien n'est pas dans le modèle de contrainte : il est dans la confiance excessive qu'il accorde au modèle d'objectif, ce qu'aucune prédiction ne peut révéler. Seule la simulation le montre.

\begin{piege}[ : mieux résoudre un mauvais modèle]
Résoudre plus exactement le problème du modèle n'aide que si le modèle est bon là où l'on va. Un solveur meilleur, quantique ou non, appliqué à un modèle qui extrapole, trouve plus sûrement ses erreurs. C'est l'argument classique des \emph{régions de confiance} en optimisation par modèle de substitution (TuRBO, par exemple) : limiter la recherche là où le modèle a des données.
\end{piege}

\begin{resultat}
\begin{itemize}
\item Sur le problème CST, avec XFOIL, le lagrangien vérifié ne bat pas les bits d'écart : 0,84 contre 0,88 d'hypervolume relatif à 300 simulations (non significatif).
\item Il respecte mieux la contrainte de $C_{L,\max}$, mais il pousse le modèle en $z$ dans ses coins, où XFOIL échoue : la pénalité à bits d'écart régularisait par accident.
\item Le classifieur réduit les échecs et aide en début de budget, sans changer le résultat final.
\item Une étude hors ligne sur des prédictions ne peut pas voir un défaut qui tient à la confiance accordée au modèle d'objectif.
\end{itemize}
\end{resultat}

\exercices
\begin{exercice}\label{ex:camp-signe}
Avec 8 graines et un test de signe bilatéral, à partir de combien de graines favorables sur 8 un écart devient-il significatif au seuil de 5~\% ? Que peut-on conclure d'un résultat « 2 graines sur 8 » ?
\end{exercice}
\begin{exercice}\label{ex:camp-coins}
Une parabole en une variable, $\hat y(z)=a z^2+bz+c$ sur $z\in[-1,1]$, ajustée sur des points tous situés dans $[-0{,}5;\,0{,}5]$. (a) Si $a<0$, où est son minimum sur $[-1,1]$ ? (b) Que vaut l'incertitude de la prédiction en $z=\pm1$ comparée à celle en $z=0$ ? (c) Expliquer pourquoi un solveur qui trouve exactement ce minimum est plus exposé qu'un solveur bruité.
\end{exercice}

\begin{pistes}
\begin{itemize}
\item Campagne \code{fm\_lagrange\_clf} (8 graines, une quarantaine de minutes) : la machine de factorisation va peu aux extrêmes, c'est l'association naturelle.
\item Une région de confiance dans le QUBO : la distance de Hamming à un profil de référence est \emph{linéaire} en bits, donc une contrainte « au plus $k$ bits changés » entre directement dans le lagrangien.
\item Régulariser le modèle en $z$ (ridge plus fort) avant de lui rendre le spectre.
\end{itemize}
\end{pistes}

